\chapter{Resultados y discusión}
\label{chap:resultados}

\drop{E}{ste} capítulo consolida los resultados globales de \dataqual una vez
completados los seis sprints descritos en el capítulo~\ref{chap:desarrollo}.
Presenta los resultados cuantitativos y cualitativos del sistema, verifica el
cumplimiento de los objetivos específicos OE1--OE5, identifica las fortalezas
y limitaciones de la plataforma, y la compara con las herramientas existentes
revisadas en el capítulo~\ref{chap:antecedentes}.


%% =========================================================================
\section{Resultados alcanzados}
\label{sec:resultados-alcanzados}
%% =========================================================================

\dataqual es una plataforma funcional de extremo a extremo que cumple con los
cinco objetivos específicos definidos en la
sección~\ref{sec:objetivos-especificos}.

\subsection{Resultados cuantitativos}

El Cuadro~\ref{tab:resultados-cuantitativos} resume las principales
magnitudes del sistema implementado.

\begin{table}[htbp]
  \centering
  \caption{Resultados cuantitativos de la implementación}
  \label{tab:resultados-cuantitativos}
  \begin{tabularx}{\textwidth}{X r}
    \toprule
    \textbf{Dimensión} & \textbf{Cantidad} \\
    \midrule
    Métricas de calidad evaluables (Quality Score) & 7     \\
    Clases de métricas implementadas (incl.\ outliers para EDA) & 8 \\
    Servicios Docker orquestados                    & 8     \\
    Rutas / páginas del \emph{frontend}             & 20    \\
    Componentes React reutilizables                 & 50+   \\
    Módulos funcionales de \ac{API} (\emph{blueprints} principales) & 7 \\
    \emph{Blueprints} Flask totales (incl.\ sub-rutas REST)        & 10 \\
    Tablas en base de datos                         & 11    \\
    Patrones de formato predefinidos (exactitud sintáctica) & 13 \\
    Tipos de \emph{issue} (por métrica y subtipo)   & 10+   \\
    Migraciones de base de datos                    & 6     \\
    Módulos de pruebas automatizadas (pytest)       & 12    \\
    \bottomrule
  \end{tabularx}
\end{table}

La batería de pruebas automatizadas se organiza en doce módulos pytest
agrupados por dominio funcional, cubriendo autenticación, \ac{API} de
proyectos, versionado de datasets a nivel de servicio y de \ac{API},
motor de métricas en sus versiones base y extendida, \emph{Quality
Score}, \emph{Quality Gates} con umbrales y veredicto,
\emph{fingerprinting}, comparación con línea base y autorización contra
accesos indebidos (\ac{IDOR}); el detalle de cada suite se describe en
la sección~\ref{sec:impl-pruebas} y el plan de pruebas integral
(estrategia, niveles, casos de prueba y trazabilidad con requisitos) se
formaliza en el Anexo~\ref{anx:pruebas}.  Las pruebas se ejecutan
con \texttt{pytest} desde el directorio \texttt{backend/} y pueden
generar informes de cobertura mediante la opción \texttt{-{}-cov}.
Estas suites unitarias y de integración se complementaron con pruebas
manuales de extremo a extremo sobre los flujos completos de la
plataforma (registro, creación de proyecto, carga de dataset,
configuración de métricas, ejecución de evaluaciones y consulta de
resultados), que verificaron la correcta integración de todos los
componentes en escenarios realistas.  La automatización de estas
pruebas \emph{end-to-end} con herramientas como Playwright o Cypress
queda como tarea pendiente y se identifica explícitamente como
limitación en la sección~\ref{sec:limitaciones}.

\subsection{Esfuerzo y desviaciones}
\label{sec:esfuerzo}

El Cuadro~\ref{tab:esfuerzo} consolida la estimación inicial de horas frente
al esfuerzo real invertido en cada sprint.

\begin{table}[htbp]
  \centering
  \caption{Esfuerzo estimado vs.\ real por sprint}
  \label{tab:esfuerzo}
  \begin{tabularx}{\textwidth}{l c c r r r}
    \toprule
    \textbf{Sprint} & \textbf{Período} & \textbf{SP} &
      \textbf{H.\ est.} & \textbf{H.\ real} & \textbf{Desv.} \\
    \midrule
    S1 — Arquitectura y setup   & oct.--nov.\ 2025 & 13 & 52  & 60  & +15,4\,\% \\
    S2 — Backend core           & dic.\ 2025--ene.\ 2026 & 18 & 72  & 85  & +18,1\,\% \\
    S3 — Motor de métricas      & feb.\ 2026       & 16 & 64  & 70  & +9,4\,\%  \\
    S4 — Frontend y dashboard   & mar.\ 2026       & 14 & 56  & 65  & +16,1\,\% \\
    S5 — Funcionalidades avanzadas & abr.\ 2026       & 10 & 40  & 55  & +37,5\,\% \\
    S6 — Pruebas y documentación & may.--jun.\ 2026 & — & 80  & 90  & +12,5\,\% \\
    \midrule
    \textbf{Total}              &                  & \textbf{71} &
      \textbf{364} & \textbf{425} & \textbf{+16,8\,\%} \\
    \bottomrule
  \end{tabularx}
  \medskip\par\footnotesize\raggedright
  \textit{Nota.}~La columna SP agrega, por sprint, las historias de
  usuario más las tareas técnicas que conformaron el plan de cada
  iteración (véase Anexo~\ref{anx:backlog}).  El backlog original
  (Tabla~\ref{tab:backlog-hu}) contabiliza 71~SP en historias de
  usuario, de los cuales se entregaron 68~SP (HU-15 fue descartada
  del alcance final); las tareas técnicas se computan adicionalmente
  como esfuerzo ejecutado.
\end{table}

La desviación global del +16,8\,\% se mantiene dentro del rango habitual para
proyectos de software individuales.  El sprint con mayor desviación fue el S5
(+37,5\,\%), causado por problemas de integración entre el sistema de
\emph{fingerprinting} y el módulo de comparación con la línea base que no
fueron detectados en la planificación.  El sprint con menor desviación fue
el S3 (+9,4\,\%), cuyo motor de métricas aprovechó la arquitectura de
\emph{plugins} establecida en S2, reduciendo el esfuerzo de integración
previsto.  La tendencia a subestimar el esfuerzo es consistente en todos los
sprints y refleja la dificultad de estimar tareas de integración entre
servicios distribuidos.

\subsection{Verificación del cumplimiento de objetivos}

El Cuadro~\ref{tab:cumplimiento} recoge, para cada objetivo específico, la
evidencia de su cumplimiento.

\begin{table}[htbp]
  \centering
  \caption{Verificación del cumplimiento de los objetivos específicos}
  \label{tab:cumplimiento}
  \begin{tabularx}{\textwidth}{l X c}
    \toprule
    \textbf{Obj.} & \textbf{Evidencia} & \textbf{Estado} \\
    \midrule
    OE1 & Módulo de ingesta en formato \ac{CSV} con parseo robusto mediante
      matriz de \emph{fallback} (codificaciones $\times$ separadores);
      almacenamiento en MinIO; extracción automática de esquema;
      versionado padre-hijo con etiquetas personalizables. & Cumplido \\
    OE2 & Siete métricas implementadas como \emph{plugins}; parametrización
      completa; sistema de pesos; plantillas reutilizables; ejecución
      asíncrona con Celery. & Cumplido \\
    OE3 & Seis \emph{blueprints} funcionales planificados, ampliados a
      siete con la adición de \texttt{patterns} (validación
      reutilizable); autenticación \ac{JWT}; formato de respuesta
      estandarizado; paginación; validación Marshmallow. & Superado \\
    OE4 & Más de 20 páginas; cuadros de mando interactivos; \ac{EDA} con
      histogramas, diagramas de caja y matrices de correlación; diseño
      adaptable. & Cumplido \\
    OE5 & Sistema de \emph{issues} con cuatro severidades;
      \emph{fingerprinting} \ac{SHA}-256; comparación con línea base;
      \emph{Quality Gates} con tres estados. & Cumplido \\
    \bottomrule
  \end{tabularx}
\end{table}


Los cinco objetivos se cumplen con el sistema tal como está desplegado.  El
único objetivo superado es OE3: la \ac{API} creció de seis a siete
\emph{blueprints} funcionales durante el desarrollo, y los \emph{blueprints}
de organización técnica elevaron el total a diez.  Los objetivos transversales
de seguridad, escalabilidad, mantenibilidad y portabilidad se cumplen; las
pruebas \emph{end-to-end} del \emph{frontend} quedan como trabajo futuro
(sección~\ref{sec:trabajo-futuro}).

%% =========================================================================
\section{Fortalezas de la plataforma}
\label{sec:fortalezas}
%% =========================================================================

A partir de los resultados consolidados en la sección anterior, las
fortalezas de \dataqual se pueden agrupar en cuatro ejes: una base
arquitectónica extensible, un conjunto de capacidades funcionales
diferenciadoras, una integración nativa de la protección de información
personal y un modelo de despliegue y operación robusto.

La principal fortaleza arquitectónica reside en el \textbf{patrón de
\emph{plugins} del motor de métricas}, materializado mediante la clase
abstracta \texttt{BaseMetric} y el diccionario
\texttt{METRIC\_REGISTRY}~\citep{design_patterns}.  Esta separación entre
la interfaz pública del motor y las implementaciones concretas permite
incorporar una nueva métrica sin modificar el código existente: basta con
crear una clase heredera, registrarla en el diccionario y, opcionalmente,
exponerla en la interfaz de configuración.  Su valor quedó demostrado
durante el propio desarrollo, en el que el catálogo creció de tres
métricas en el Sprint~3 a siete al cierre del Sprint~5 sin necesidad de
reescribir el orquestador de evaluaciones.  La misma filosofía de
extensibilidad se traslada al \emph{frontend}, organizado en torno a más
de cincuenta componentes React reutilizables que permiten ensamblar
nuevas pantallas con bajo coste marginal.

Sobre esa base arquitectónica, \dataqual incorpora cuatro capacidades
funcionales que la diferencian de las herramientas analizadas en el
capítulo~\ref{chap:antecedentes}.  El \textbf{sistema de
\emph{fingerprinting}}, basado en un \emph{hash} \ac{SHA}-256 truncado y
determinista por \emph{issue}, permite trazar la evolución de un mismo
problema entre ejecuciones sucesivas y clasificar los \emph{issues}
automáticamente como nuevos, recurrentes o corregidos: una capacidad
poco común en herramientas de código abierto, que típicamente reportan
listados planos sin continuidad temporal.  Los \textbf{\emph{Quality
Gates}} trasladan al dominio de los datos un concepto consolidado en la
calidad del código~\citep{sonarqube}, ofreciendo un veredicto
PASSED/WARNING/FAILED automatizable que puede integrarse en
\emph{pipelines} de \ac{MLOps} como criterio de promoción de datasets
entre entornos.  El \textbf{módulo de perfilado interactivo} concentra
en una sola pantalla histogramas, diagramas de caja, matrices de
correlación de Pearson y Spearman y detección de \emph{outliers} con
factor \ac{IQR} configurable, sin obligar al usuario a abandonar la
plataforma.  Por último, las \textbf{métricas específicas para
\emph{machine learning}} (equilibrio de clases basado en entropía de
Shannon y diversidad del \emph{dataset}) responden a necesidades del
ciclo de vida de modelos que las herramientas generalistas de calidad
de datos no cubren.

La \textbf{protección nativa de información personal identificable
(\ac{PII})} es la fortaleza con mayor coste de implementación pero
también la de mayor valor diferencial en contextos regulados por el
\ac{RGPD}.  Lejos de tratarse de una capa filtro añadida \emph{a
posteriori}, el enmascaramiento se ha implementado como un contrato que
toda métrica debe respetar: las clases heredan de \texttt{BaseMetric}
un método de enmascaramiento que se invoca antes de cualquier escritura
de un valor a logs, \emph{issues} o respuestas de la \ac{API},
garantizando que ninguna columna marcada como sensible filtre valores a
través de cualquier camino del sistema.  Esta integración estructural
marca la diferencia entre «la herramienta puede ofuscar \ac{PII} si se
configura» y «la herramienta no puede revelar \ac{PII} por
construcción».

Desde el punto de vista operativo, la plataforma se distingue por ser
\textbf{íntegramente autoalojada y desplegable con un único comando}
(\texttt{docker-compose~up}): los ocho servicios se levantan en una red
interna, sin dependencia de servicios en la nube propietarios y
manteniendo los datos del usuario bajo su control total.  La capa de
ingesta incorpora un \textbf{parseo robusto de \ac{CSV}} mediante una
matriz de \emph{fallback} que combina cuatro codificaciones, seis
estrategias de separador y los dos motores de Pandas, capaz de gestionar
ficheros reales con BOM, codificaciones mixtas o delimitadores poco
habituales que provocan errores en analizadores estándar.  El
\textbf{procesamiento asíncrono con Celery}~\citep{celery} se complementa
con un \emph{watchdog} que detecta evaluaciones atascadas y un servicio
híbrido que recurre al modo síncrono ante indisponibilidad del
\emph{worker}, asegurando que las evaluaciones se completen incluso en
escenarios degradados.


%% =========================================================================
\section{Limitaciones}
\label{sec:limitaciones}
%% =========================================================================

A pesar de los resultados alcanzados, la plataforma presenta limitaciones
que conviene reconocer explícitamente.  Estas se pueden agrupar en cuatro
ámbitos: la escalabilidad y el rendimiento del motor de procesamiento,
la madurez de los servicios de plataforma de cara a entornos
corporativos, la cobertura de las pruebas automatizadas y el alcance de
la validación experimental.

El primer grupo afecta a la \textbf{escalabilidad y al rendimiento} bajo
cargas elevadas y constituye la restricción técnica más significativa de
la plataforma.  El motor de métricas opera sobre \emph{DataFrames} de
Pandas cargados íntegramente en memoria, lo que impone como cota
superior el tamaño de la RAM disponible y descarta procesar
\emph{datasets} de varios GB sin migrar a un motor distribuido como Dask
o Spark.  La capa de procesamiento, además, no se ha validado con
múltiples \emph{workers} de Celery en paralelo, por lo que la
escalabilidad horizontal que el diseño permite es todavía teórica.  El
\emph{frontend}, por su parte, obtiene el estado de las evaluaciones
mediante \emph{polling} periódico en lugar de notificaciones \emph{push}
por \emph{WebSocket}: aunque el patrón es funcionalmente correcto y
simplifica la arquitectura, introduce latencia perceptible y carga
innecesaria contra la \ac{API} cuando se ejecutan muchas evaluaciones
simultáneamente.  Por último, la detección de valores atípicos en el
módulo de perfilado se limita al método estadístico clásico del
\ac{IQR}; algoritmos basados en aprendizaje automático como
\emph{Isolation Forest} o DBSCAN, que detectarían anomalías
multidimensionales más sutiles, quedan fuera del alcance de esta versión.

Un segundo grupo afecta a la \textbf{madurez de los servicios de
plataforma} de cara a entornos corporativos.  La autenticación se limita
a \ac{JWT} con usuario y contraseña: no se han implementado OAuth~2.0,
\emph{Single Sign-On} corporativo ni autenticación multifactor, por lo
que la plataforma no está preparada todavía para integrarse en entornos
empresariales con identidad federada.  Los resultados, además, no pueden
exportarse a formatos de informe formal como PDF o Excel, lo que
dificulta el reporte a perfiles ajenos al equipo técnico y la
incorporación de los hallazgos a documentación de auditoría.

El tercer ámbito es la cobertura de \textbf{pruebas automatizadas}, que
aunque amplia en componentes individuales (motor de métricas,
\emph{Quality Score}, \emph{Quality Gates}, \emph{fingerprinting},
comparación con línea base, autorización \ac{IDOR} y versionado de
\emph{datasets}), no incluye pruebas \emph{end-to-end} que validen el
flujo completo desde la subida de un \emph{dataset} hasta la emisión del
veredicto del \emph{Quality Gate}.  Esta validación se realizó de forma
manual durante el desarrollo, pero su automatización con herramientas
como Playwright o Cypress queda como tarea pendiente y constituye la
mayor brecha en la estrategia de aseguramiento de calidad del propio
proyecto.

Finalmente, conviene reconocer también una limitación de carácter
metodológico relativa al \textbf{alcance de la validación
experimental}.  La verificación funcional del sistema se ha realizado
por el propio autor, complementada con un conjunto amplio de pruebas
automatizadas, pero no se ha llevado a cabo una validación con usuarios
externos del dominio de calidad de datos o de ingeniería de \ac{ML} que
permita evaluar empíricamente la usabilidad de la interfaz, la
adecuación del catálogo de métricas a casos de uso reales o el ajuste
de los umbrales por defecto de los \emph{Quality Gates}.  La
validación realizada acredita la corrección técnica del sistema, pero
no su idoneidad operativa para perfiles externos al desarrollo.


%% =========================================================================
\section{Comparación con herramientas existentes}
\label{sec:comparacion-final}
%% =========================================================================

El Cuadro~\ref{tab:comparacion-final} retoma la comparación del
capítulo~\ref{chap:antecedentes} enriquecida con perspectiva de
implementación: no solo si la funcionalidad existe, sino qué implicó
reproducirla o superarla en la práctica.

\begin{table}[htbp]
  \centering
  \caption{Comparación de \dataqual con herramientas existentes (perspectiva post-implementación)}
  \label{tab:comparacion-final}
  \begin{tabularx}{\textwidth}{X c c c c c}
    \toprule
    \textbf{Criterio} & \textbf{\dataqual} &
      \textbf{Great Exp.} & \textbf{Deequ} & \textbf{DQOps} & \textbf{Soda} \\
    \midrule
    \ac{UI} integrada          & Sí  & No       & No         & Sí      & Parcial$^a$ \\
    Requiere código            & No  & Sí       & Sí         & Parcial & No \\
    \emph{Quality Gates}       & Sí  & No       & No         & Sí      & Sí \\
    \emph{Fingerprinting}      & Sí  & No       & No         & Parcial & No \\
    Diff entre ejecuciones     & Sí  & No       & Parcial    & Sí      & Sí \\
    Métricas para \ac{ML}      & Sí  & No       & Parcial    & No      & No \\
    Enmascaramiento \ac{PII}   & Sí  & No       & No         & No      & No \\
    Versionado de datasets     & Sí  & No       & No         & No      & No \\
    Autoalojado                & Sí  & Sí       & Sí         & Sí      & Sí \\
    Parseo robusto de archivos & Sí  & Parcial  & No         & No      & Parcial \\
    Madurez                    & Prototipo & Producción & Producción & Producción & Producción \\
    \bottomrule
  \end{tabularx}
  \medskip\par\footnotesize\raggedright
  \textsuperscript{a}~La interfaz web de Soda (Soda Cloud) está disponible únicamente en el nivel comercial.
\end{table}

La implementación aportó perspectivas que no eran evidentes en el análisis
previo del estado del arte.  El \textbf{parseo robusto de archivos} resultó
ser un diferenciador real: en las pruebas con conjuntos de datos reales, una
fracción significativa de los ficheros \ac{CSV} presentó combinaciones de
codificación y delimitador no estándar que ninguna herramienta \emph{code-first}
gestiona automáticamente sin configuración explícita.  El
\textbf{enmascaramiento nativo de \ac{PII}} supuso un coste de implementación
considerable~\textemdash todas las clases de métricas deben respetar la lista
de columnas sensibles~\textemdash pero su valor es especialmente alto en
contextos regulados por el \ac{RGPD}, donde ninguna de las alternativas
analizadas lo ofrece de forma integrada.

\dataqual no pretende competir con herramientas maduras de producción como
Great Expectations~\citep{greatexpectations} en escala o número de métricas
predefinidas, pero ocupa un nicho diferenciado al combinar interfaz web
completa, métricas específicas para \ac{ML}, \emph{fingerprinting} determinista,
\emph{Quality Gates} y protección nativa de \ac{PII} en un paquete autoalojado
que no requiere escribir código.


% Local Variables:
%  coding: utf-8
%  mode: latex
%  mode: flyspell
%  ispell-local-dictionary: "castellano8"
% End:
